%%%%%%%%%%%%%%%%%%%%%%%%%%%%%%%%%%%%%%%%%%%%%%%%%%%%%%%%%%%%%%%%%%%%%%
%% Shared pieces for the l04_afe Gm-C cell figures.
%%
%% l4_gmc, l4_gmc_diff, l4_gmc_diff1 and l4_gmc1st all draw the
%% transconductor the same way the original artwork does: a trapezoid with a
%% tall input edge carrying the two input terminals and a short output edge.
%% They are consecutive slides making the same point four ways, so they have
%% to be on one geometry.
%%
%% l4_gmc_diff and l4_gmc_diff1 differ only in the output terminal marks and
%% in which way the two G_m V_i arrows point -- that difference is the whole
%% slide -- so \gmcDiffFrame draws everything else and leaves those to the
%% caller, the same way sc_lib's \scIntroFrame does.
%%
%% Every name is prefixed and every macro name is letters only: TeX reads
%% \gmcC1 as \gmcC followed by a stray 1, and fails at use rather than at
%% definition.  Computed coordinates are wrapped in braces, because TikZ
%% scans a coordinate for its closing parenthesis and a '(' inside the
%% expression ends it early.
%%%%%%%%%%%%%%%%%%%%%%%%%%%%%%%%%%%%%%%%%%%%%%%%%%%%%%%%%%%%%%%%%%%%%%

\newcommand{\gmcW}{1.5}    % body width
\newcommand{\gmcHi}{1.5}   % half height of the tall input edge
\newcommand{\gmcHo}{0.8}   % half height of the output edge, which is also
                           % where the input terminals sit

% The body, pointing right.  (#1,#2) is the centre of the input edge.
\newcommand{\gmcBody}[3]{%
  \draw (#1,{#2-\gmcHi}) -- (#1,{#2+\gmcHi})
     -- ({#1+\gmcW},{#2+\gmcHo}) -- ({#1+\gmcW},{#2-\gmcHo}) -- cycle;
  \node at ({#1+\gmcW/2},#2) {#3};
}

% Mirrored: body to the left of the input edge, output on the left.
\newcommand{\gmcBodyL}[3]{%
  \draw (#1,{#2-\gmcHi}) -- (#1,{#2+\gmcHi})
     -- ({#1-\gmcW},{#2+\gmcHo}) -- ({#1-\gmcW},{#2-\gmcHo}) -- cycle;
  \node at ({#1-\gmcW/2},#2) {#3};
}

% Terminal marks, drawn as strokes so they scale with the line work.
\newcommand{\gmcPlus}[2]{%
  \draw ({#1-0.17},#2) -- ({#1+0.17},#2);
  \draw (#1,{#2-0.17}) -- (#1,{#2+0.17});
}
\newcommand{\gmcMinus}[2]{\draw ({#1-0.17},#2) -- ({#1+0.17},#2);}

% A capacitor in a vertical run from (#1,#2) down to (#1,#3).
\newcommand{\gmcVcap}[3]{%
  \draw (#1,#2) -- (#1,{(#2+#3)/2+0.05});
  \draw ({#1-0.2},{(#2+#3)/2+0.05}) -- ({#1+0.2},{(#2+#3)/2+0.05});
  \draw ({#1-0.2},{(#2+#3)/2-0.05}) -- ({#1+0.2},{(#2+#3)/2-0.05});
  \draw (#1,{(#2+#3)/2-0.05}) -- (#1,#3);
}

% A capacitor in a horizontal run from (#1,#2) to (#3,#2).
\newcommand{\gmcHcap}[3]{%
  \draw (#1,#2) -- ({(#1+#3)/2-0.05},#2);
  \draw ({(#1+#3)/2-0.05},{#2-0.2}) -- ({(#1+#3)/2-0.05},{#2+0.2});
  \draw ({(#1+#3)/2+0.05},{#2-0.2}) -- ({(#1+#3)/2+0.05},{#2+0.2});
  \draw ({(#1+#3)/2+0.05},#2) -- (#3,#2);
}

% ---------------------------------------------------------------------------
% Everything l4_gmc_diff and l4_gmc_diff1 share.  Body at (1.2,2.8), rails at
% 3.6 and 2.0, C across the outputs at x=4.3.  The caller adds the two output
% terminal marks at (2.35,3.6) and (2.35,2.0) and the two current arrows.
% ---------------------------------------------------------------------------
\newcommand{\gmcDiffFrame}{%
  \gmcBody{1.2}{2.8}{$G_{m}$}
  \draw (0.3,3.6) -- (1.2,3.6);
  \draw (0.3,2.0) -- (1.2,2.0);
  \gmcPlus{1.55}{3.6}
  \gmcMinus{1.55}{2.0}
  \node[anchor=east] at (0.25,2.8) {$V_{i}$};
  \node at (0.55,3.32) {$+$};
  \node at (0.55,2.28) {$-$};

  \draw (2.7,3.6) -- (4.3,3.6);
  \draw (2.7,2.0) -- (4.3,2.0);
  \gmcVcap{4.3}{3.6}{2.0}
  \node[anchor=west] at (4.45,3.42) {$+\ V_{o}$};
  \node[anchor=west] at (4.45,2.18) {$-$};
  \node[anchor=west] at (4.72,2.8) {$C$};
}
